%% PATRÓN · Flujo con cajas y bucle (un paso de un algoritmo, una iteración).
\providecolor{textochip}{named}{white} % por si la charla tiene un estilo/ anterior a 1.1.0
%% Cajas con relleno suave del color de su familia y BORDE NEGRO (tinta, 0.6pt); un marco gris claro
%% agrupa lo que se repite y una flecha de color cierra el bucle con su condición. Entrada y salida
%% fuera del marco (n -> n+1). Original: charla01/algoritmo_imp.tex.
%% Estilos con #1: este fichero va con \input, nunca pegado dentro de un frame.
\begin{tikzpicture}[
    caja/.style={fill=#1!12!fondo,draw=tinta,line width=0.6pt,rounded corners=3pt,align=left,inner sep=5pt,
      font=\scriptsize,text width=2.5cm,minimum height=2.4cm},
    sol/.style={fill=#1,text=textochip,rounded corners=1.5pt,inner sep=1.4pt,font=\tiny\bfseries,
      minimum height=0pt,minimum width=0pt,text width={},align=center},   % TikZ anidado: resetear
    flecha/.style={-{Stealth[length=5pt]},tinta,line width=0.8pt}]
  \newcommand\etiqueta[2]{\tikz[baseline=(etq@.base)]\node[sol=#1](etq@){#2};} % nombre raro: los nodos son globales
  \node[caja=serieA] (a) at (0,0)
    {\textbf{1 · Ensamblar}\\[3pt] matriz y residuo\\[3pt] {\color{tenue}lo que cambia por iteración}};
  \node[caja=serieC,right=0.4cm of a,text width=3.4cm] (b)
    {\textbf{2 · Resolver} \etiqueta{serieC}{SOLVER}\\[3pt] $A\,\boldsymbol x = \boldsymbol b$\\[3pt]
     {\color{tenue}GMRES + precondicionador}};
  \node[caja=serieD,right=0.4cm of b] (c)
    {\textbf{3 · Actualizar}\\[3pt] $\boldsymbol x^{k+1} = \boldsymbol x^k + \delta\boldsymbol x$};
  \draw[flecha] (a) -- (b); \draw[flecha] (b) -- (c);
  \begin{scope}[on background layer]
    \node[fill=tinta!4!fondo,rounded corners=5pt,inner xsep=8pt,inner ysep=10pt,fit=(a)(b)(c)] (it) {};
  \end{scope}
  \node[anchor=south west,font=\scriptsize\bfseries,inner sep=1pt] at ($(it.north west)+(0.15,0.02)$)
    {Iteración no lineal};
  \draw[flecha,acento] (c.south) -- ++(0,-0.62) -| node[pos=0.25,below,font=\scriptsize,text=tinta]
    {¿residuo $> 10^{-3}$? otra iteración} (a.south);
  \draw[flecha] ($(it.west)+(-0.65,0)$) node[left,font=\small]{$n$} -- (it.west);
  \draw[flecha] (it.east) -- ++(0.55,0) node[right,font=\small]{$n{+}1$};
\end{tikzpicture}
